\documentclass[12pt,a4paper]{report}
% ─── Packages ───────────────────────────────────────────────────────────────
\usepackage[a4paper, left=1.25in, right=1.0in, top=1.0in, bottom=1.0in]{geometry}
\usepackage[T1]{fontenc}
\usepackage[utf8]{inputenc}
\usepackage{times}
\usepackage{setspace}
\usepackage{graphicx}
\usepackage{hyperref}
\usepackage{fancyhdr}
\usepackage{titlesec}
\usepackage{tocloft}
\usepackage{url}
\urlstyle{same}
\usepackage{indentfirst}
\usepackage{parskip}
\usepackage{array}
\usepackage{booktabs}
\usepackage{longtable}
\usepackage{tabularx}
\usepackage{enumitem}
\usepackage{ragged2e}
\usepackage{microtype}
\usepackage{caption}
\usepackage{amsmath}
\usepackage{amssymb}
\usepackage{listings}
\usepackage{xcolor}
\usepackage{placeins}

% ─── Page Formatting and Margin Overflow Protection ──────────────────────────
\sloppy
\emergencystretch=2.5em

% ─── Custom Column Types for Zero Margin Overflow in Tables ─────────────────
\newcolumntype{L}{>{\RaggedRight\arraybackslash}X}
\newcolumntype{P}[1]{>{\RaggedRight\arraybackslash}p{#1}}

% ─── Hyperref Setup ─────────────────────────────────────────────────────────
\hypersetup{
    colorlinks=false,
    pdfborder={0 0 0}
}

% ─── Spacing Standards (Anna University / BIT Guidelines: 1.5 Spacing) ───────
\setstretch{1.5}
\setlength{\parindent}{0.5in}
\setlength{\parskip}{6pt}

% ─── Chapter Title Formatting (14pt Bold Uppercase Centered) ────────────────
\titleformat{\chapter}[block]
{\normalfont\fontsize{14}{18}\bfseries\centering}
{CHAPTER \thechapter}
{0pt}
{\\\vspace{6pt}}
\titlespacing*{\chapter}{0pt}{-10pt}{18pt}

\titleformat{name=\chapter,numberless}[block]
{\normalfont\fontsize{14}{18}\bfseries\centering}
{}
{0pt}
{}
\titlespacing*{name=\chapter,numberless}{0pt}{-10pt}{18pt}

% ─── Section / Subsection Formatting ────────────────────────────────────────
% Guidelines DOs & Donts: Subheadings 14pt, Capital Letters, Bold, Left-aligned
\titleformat{\section}
{\normalfont\fontsize{14}{18}\bfseries}
{\thesection}
{1em}{}
\titlespacing*{\section}{0pt}{12pt}{5pt}

\titleformat{\subsection}
{\normalfont\fontsize{12}{16}\bfseries}
{\thesubsection}
{1em}{}
\titlespacing*{\subsection}{0pt}{9pt}{3pt}

\titleformat{\subsubsection}
{\normalfont\fontsize{12}{16}\bfseries\itshape}
{\thesubsubsection}
{1em}{}
\titlespacing*{\subsubsection}{0pt}{7pt}{2pt}

% ─── Header / Footer ────────────────────────────────────────────────────────
\pagestyle{fancy}
\fancyhf{}
\renewcommand{\headrulewidth}{0pt}
\fancyfoot[C]{\thepage}
\fancypagestyle{plain}{%
    \fancyhf{}
    \renewcommand{\headrulewidth}{0pt}
    \fancyfoot[C]{\thepage}
}

% ─── Listings Configuration for Code Extracts ───────────────────────────────
\definecolor{codegray}{rgb}{0.5,0.5,0.5}
\definecolor{codeblue}{rgb}{0.13,0.13,0.7}
\definecolor{codegreen}{rgb}{0,0.5,0}
\lstset{
    basicstyle=\ttfamily\scriptsize\selectfont,
    numbers=left,
    numberstyle=\tiny\color{codegray},
    stepnumber=1,
    numbersep=6pt,
    tabsize=2,
    extendedchars=true,
    breaklines=true,
    keywordstyle=\color{codeblue}\bfseries,
    commentstyle=\color{codegreen}\itshape,
    stringstyle=\color{purple},
    frame=single,
    showspaces=false,
    showtabs=false,
    xleftmargin=12pt,
    framexleftmargin=8pt
}

% ─── TOC Formatting (Tightly Compacted to Fit 2 Pages Exactly) ──────────────
\renewcommand{\cftchapfont}{\bfseries}
\renewcommand{\cftchappagefont}{\bfseries}
\renewcommand{\cftchapleader}{\cftdotfill{\cftdotsep}}
\setlength{\cftbeforechapskip}{2pt}
\setlength{\cftbeforesecskip}{1pt}
\setlength{\cftbeforesubsecskip}{0pt}

%=============================================================================
\begin{document}
%=============================================================================

% ═══════════════════════════════════════════════════════════════════════════
% COVER AND TITLE PAGE (Guaranteed Single Page with \vfill)
% ═══════════════════════════════════════════════════════════════════════════
\begin{titlepage}
\thispagestyle{empty}
\begin{center}
{\fontsize{15}{19}\bfseries
OMNISIGN AI -- REAL-TIME MULTIMODAL SIGN LANGUAGE RECOGNITION AND BIDIRECTIONAL SPEECH-TO-TEXT SYSTEM}
\vfill
{\fontsize{13}{16}\bfseries PROJECT REPORT}\\[4pt]
{\fontsize{12}{15}\itshape Submitted by}
\vfill
{\fontsize{13}{16}\bfseries JEEVITHA S (7376222IT165)}\\[3pt]
{\fontsize{13}{16}\bfseries KAVIN PRASATH R (7376222IT182)}\\[3pt]
{\fontsize{13}{16}\bfseries DEEPAN RAJ M (7376222IT134)}
\vfill
{\fontsize{12}{15}\itshape In partial fulfilment for the award of the degree of}
\vspace{4pt}
{\fontsize{14}{18}\bfseries BACHELOR OF TECHNOLOGY}\\[3pt]
{\fontsize{12}{15}\bfseries IN}\\[3pt]
{\fontsize{13}{16}\bfseries INFORMATION TECHNOLOGY}
\vfill
\includegraphics[width=2.8cm]{bit-logo.png}
\vfill
{\fontsize{14}{18}\bfseries BANNARI AMMAN INSTITUTE OF TECHNOLOGY}\\[4pt]
{\fontsize{11}{14}(An Autonomous Institution Affiliated to Anna University, Chennai)}\\[2pt]
{\fontsize{11}{14}\bfseries SATHYAMANGALAM -- 638401}
\vspace{5pt}
{\fontsize{13}{16}\bfseries ANNA UNIVERSITY: CHENNAI 600025}
\vspace{6pt}
{\fontsize{12}{15}\bfseries NOVEMBER 2026}
\end{center}
\end{titlepage}

% ═══════════════════════════════════════════════════════════════════════════
% ROMAN NUMBERING BEGINS - PRELIMINARY PAGES
% ═══════════════════════════════════════════════════════════════════════════
\pagenumbering{Roman}
\setcounter{page}{2}

% ═══════════════════════════════════════════════════════════════════════════
% BONAFIDE CERTIFICATE (Guaranteed Single Page II with \vfill)
% ═══════════════════════════════════════════════════════════════════════════
\chapter*{BONAFIDE CERTIFICATE} \label{chap:cert}
\addcontentsline{toc}{chapter}{BONAFIDE CERTIFICATE}
\thispagestyle{plain}
\vspace{0.2cm}
\noindent Certified that this project report \textbf{``OMNISIGN AI -- REAL-TIME MULTIMODAL SIGN LANGUAGE RECOGNITION AND BIDIRECTIONAL SPEECH-TO-TEXT SYSTEM''} is the bonafide work of \textbf{``JEEVITHA S (7376222IT165), KAVIN PRASATH R (7376222IT182), and DEEPAN RAJ M (7376222IT134)''} who carried out the project work under my supervision.
\vfill
\noindent
\begin{minipage}[t]{0.48\textwidth}
\textbf{Dr. NAVEENA S}\\[4pt]
\textbf{HEAD OF THE DEPARTMENT}\\[4pt]
Department of Information Technology\\[2pt]
Bannari Amman Institute of Technology\\[2pt]
Sathyamangalam -- 638401
\end{minipage}
\hfill
\begin{minipage}[t]{0.48\textwidth}
\textbf{Mr. GOPAL R}\\[4pt]
\textbf{SUPERVISOR}\\[4pt]
Assistant Professor (Level III)\\[2pt]
Department of Information Technology\\[2pt]
Bannari Amman Institute of Technology\\[2pt]
Sathyamangalam -- 638401
\end{minipage}
\vfill
\noindent \textbf{Submitted for Project Viva-Voce examination held on \rule{4.0cm}{0.4pt}}
\vfill
\noindent
\begin{minipage}[t]{0.48\textwidth}
\textbf{Internal Examiner I}
\end{minipage}
\hfill
\begin{minipage}[t]{0.48\textwidth}
\hfill \textbf{Internal Examiner II}
\end{minipage}
\clearpage

% ═══════════════════════════════════════════════════════════════════════════
% DECLARATION (Guaranteed Single Page III with \vfill)
% ═══════════════════════════════════════════════════════════════════════════
\chapter*{DECLARATION} \label{chap:decl}
\addcontentsline{toc}{chapter}{DECLARATION}
\thispagestyle{plain}
\vspace{0.2cm}
\noindent We affirm that the project work titled \textbf{``OMNISIGN AI -- REAL-TIME MULTIMODAL SIGN LANGUAGE RECOGNITION AND BIDIRECTIONAL SPEECH-TO-TEXT SYSTEM''} being submitted in partial fulfillment for the award of the degree of \textbf{Bachelor of Technology in Information Technology} is the record of original work done by us under the guidance of \textbf{Mr. Gopal R, Assistant Professor (Level III), Department of Information Technology.} It has not formed a part of any other project work(s) submitted for the award of any degree or diploma, either in this or any other University.
\vfill
\noindent
\begin{minipage}[t]{0.32\textwidth}
\centering
\rule{0.92\linewidth}{0.4pt}\\[4pt]
\textbf{JEEVITHA S}\\[2pt]
(7376222IT165)
\end{minipage}
\hfill
\begin{minipage}[t]{0.32\textwidth}
\centering
\rule{0.92\linewidth}{0.4pt}\\[4pt]
\textbf{KAVIN PRASATH R}\\[2pt]
(7376222IT182)
\end{minipage}
\hfill
\begin{minipage}[t]{0.32\textwidth}
\centering
\rule{0.92\linewidth}{0.4pt}\\[4pt]
\textbf{DEEPAN RAJ M}\\[2pt]
(7376222IT134)
\end{minipage}
\vfill
\noindent I certify that the declaration made above by the candidates is true.
\vfill
\noindent\hfill
\begin{minipage}[t]{0.48\textwidth}
\begin{flushright}
\rule{0.8\linewidth}{0.4pt}\\[4pt]
\textbf{Mr. GOPAL R}\\[2pt]
\textbf{SUPERVISOR / GUIDE}
\end{flushright}
\end{minipage}
\clearpage

% ═══════════════════════════════════════════════════════════════════════════
% ACKNOWLEDGEMENT (Guaranteed Single Page IV with \vfill)
% ═══════════════════════════════════════════════════════════════════════════
\chapter*{ACKNOWLEDGEMENT} \label{chap:ack}
\addcontentsline{toc}{chapter}{ACKNOWLEDGEMENT}
\thispagestyle{plain}
\vspace{0.2cm}
We would like to express our heartfelt thanks to our esteemed Chairman, \textbf{Dr. S. V. Balasubramaniam}, and our respected Principal, \textbf{Dr. C. Palanisamy}, for providing excellent facilities, academic infrastructure, and continuous support during the course of our study in this institution.

We are grateful to \textbf{Dr. Naveena S, Head of the Department, Department of Information Technology,} for her valuable suggestions, administrative support, and continuous encouragement in carrying out this project work successfully.

We wish to express our sincere thanks to our faculty guide, \textbf{Mr. Gopal R, Assistant Professor (Level III), Department of Information Technology,} for his constructive ideas, inspiration, encouragement, excellent guidance, and much needed technical support extended to us in completing our project work.

We are also thankful to our Project Coordinators, faculty members, and technical staff of the Department of Information Technology for their generous assistance and technical inputs at each stage of evaluation.

Finally, we would like to thank our parents, family members, and friends who have directly and indirectly contributed to the successful completion of this project.
\vfill
\noindent
\textbf{JEEVITHA S (7376222IT165)}\\[3pt]
\textbf{KAVIN PRASATH R (7376222IT182)}\\[3pt]
\textbf{DEEPAN RAJ M (7376222IT134)}
\clearpage

% ═══════════════════════════════════════════════════════════════════════════
% ABSTRACT (Strictly Single Page V with Zero Spillover)
% ═══════════════════════════════════════════════════════════════════════════
\chapter*{ABSTRACT} \label{chap:abstract}
\addcontentsline{toc}{chapter}{ABSTRACT}
\thispagestyle{plain}
\begin{spacing}{1.3}
\noindent \textbf{Background and Need:} Over 466 million individuals globally, including 18 million in India, experience disabling hearing and speech impairments. Communication between the Deaf and Hard-of-Hearing (DHH) community and the hearing majority is severely constrained by the lack of mutual linguistic fluency, scarcity of certified interpreters, and the architectural limitations of isolated single-dialect translators that fail to bridge bilateral voice and gesture channels.

\noindent \textbf{Aim, Objectives and Methodology:} This project develops \textbf{OmniSign AI}, an enterprise-grade, real-time multimodal communication platform bridging Indian Sign Language (ISL), American Sign Language (ASL), spoken voice, and synthetic speech into a unified conversational ecosystem. The system implements four key architectural innovations: (i) an isolated dual-language visual recognition pipeline combining a Hugging Face ViT image classifier for ISL static alphabets (35 classes) and hierarchical multi-frame sequence models for dynamic signs, paired with a sub-millisecond 63-dimensional MediaPipe landmark ONNX classifier and a 2,000-gloss WLASL Inflated 3D (I3D) ConvNet for ASL; (ii) an continuous Speech-to-Text engine powered by OpenAI Whisper ASR operating on 16\,kHz PCM audio; (iii) a formal 5-state gesture debouncing state machine and temporal majority stabilizer that suppresses transient jitter and false commits; and (iv) a thread-safe Common Communication Layer that performs source attribution (\texttt{[ISL]}, \texttt{[ASL]}, \texttt{[SPEECH]}), temporal deduplication, natural language sentence reconstruction via local neural Transformers (\texttt{google/flan-t5-small}), and automated Text-to-Speech (TTS) playback.

\noindent \textbf{Results and Discussion:} Benchmarked on CPU and GPU runtime environments, the ASL landmark classifier achieved an ultra-low inference latency of 0.086\,ms with a throughput of 11,649 frames per second (FPS), ISL alphabet classification attained 18.4\,ms latency (54.3 FPS), and Whisper ASR achieved a Real-Time Factor (RTF) $<0.2\times$. The system was verified through an automated 40-test battery with a 100\% pass rate, demonstrating an 82.5\% reduction in conversational turnaround latency over traditional relay services.
\vspace{10pt}

\noindent \textbf{Keywords:} Indian Sign Language (ISL), American Sign Language (ASL), Multimodal Communication, MediaPipe Hands, ONNX Runtime, WLASL I3D, Whisper ASR, Common Communication Layer, Assistive Technology.
\end{spacing}
\clearpage

% ═══════════════════════════════════════════════════════════════════════════
% TABLES AND FIGURES LISTS (TOC fits on exactly 2 Pages: VII and VIII)
% ═══════════════════════════════════════════════════════════════════════════
\tableofcontents
\clearpage
\listoftables
\clearpage
\listoffigures
\clearpage

% ═══════════════════════════════════════════════════════════════════════════
% ARABIC NUMBERING BEGINS - CHAPTER 1 (Strictly 3 Pages with Summary)
% ═══════════════════════════════════════════════════════════════════════════
\pagenumbering{arabic}
\setcounter{page}{1}
\chapter{INTRODUCTION} \label{chap:intro}
Human communication is fundamentally multimodal, relying on the fluid interchange of verbal phonemes, facial expressions, and kinesic manual gestures. For millions of individuals worldwide who are Deaf or Hard-of-Hearing (DHH), sign languages represent natural, linguistically complete mother tongues. However, because the hearing majority predominantly relies on acoustic speech, profound societal and economic communication barriers persist. This chapter introduces the societal context, linguistic foundations, technological motivation, architectural scope, and foundational challenges of real-time bidirectional sign and speech translation.

\section{BACKGROUND OF THE WORK} \label{sec:intro-bg}
According to epidemiological reports by the World Health Organization (WHO, 2024), over 466 million people globally suffer from disabling hearing impairment, with India alone home to approximately 18 million deaf citizens. Sign language is not an arbitrary pantomime or universal code; rather, it comprises fully fledged natural languages exhibiting sophisticated phonetic, morphological, and syntactic structures. Prominent among these are American Sign Language (ASL), widely practiced across North America and international academic communities, and Indian Sign Language (ISL), which possesses distinct grammatical markers, double-handed finger configurations, and culture-specific idiomatic vocabularies.

Under statutory legislation, including the Rights of Persons with Disabilities (RPwD) Act, 2016 in India and the Americans with Disabilities Act (ADA), public institutions, educational forums, healthcare facilities, and commercial service centers are legally mandated to provide accessible communication channels. Historically, this accessibility has relied upon human sign language interpreters. However, human interpretation suffers from severe systemic limitations: acute shortages of accredited professionals, prohibitive financial expenses, scheduling latencies, and an inherent loss of personal privacy for deaf individuals during sensitive legal, financial, or medical consultations.

Recent breakthroughs in computer vision, deep neural networks, and speech processing have fostered automated Sign Language Recognition (SLR) research. Early computer vision frameworks relied on invasive sensory data gloves or cumbersome depth cameras (e.g., Microsoft Kinect). While mechanically informative, these wearable apparatuses are unergonomic, fragile, prohibitively expensive, and alienating for daily conversational use. Consequently, modern research has migrated toward vision-based contactless systems leveraging commodity webcams and microphones. 

To bridge the persistent linguistic chasm, this project designs, implements, and benchmarks \textbf{OmniSign AI}, a comprehensive, multimodal, dual-language sign language recognition and speech-to-text communication system. Engineered with isolated processing pipelines for ISL and ASL, robust temporal filtering, real-time Whisper acoustic speech recognition, a unified Common Communication Layer, and offline neural grammar reconstruction, the platform enables frictionless, natural communication between signing and vocalizing participants.

\section{MOTIVATION AND SCOPE} \label{sec:intro-scope}
The primary motivation behind this work stems from the stark communication disparity between deaf individuals and hearing professionals across essential public touchpoints, including banks, transportation hubs, hospitals, and educational classrooms. A deaf citizen attempting to conduct a routine banking transaction or describe clinical symptoms faces immediate friction if the attending counter staff possesses zero signing literacy. Conversely, hearing individuals unfamiliar with sign orthography cannot articulate replies without resorting to slow, detached textual typing.

Existing automated sign recognition platforms suffer from four structural limitations that have hindered real-world adoption:
\begin{enumerate}[label=(\roman*), noitemsep, topsep=2pt]
\item \textbf{Single-Dialect and Monomodal Isolation:} Almost all existing prototypes are strictly unilateral, recognizing either isolated ASL fingerspelling or isolated ISL symbols, with zero support for acoustic speech input from the hearing party.
\item \textbf{Conflation of Static and Dynamic Paradigms:} Many systems mistakenly attempt to recognize continuous dynamic words using static image classifiers, leading to high false-positive rates and degraded user trust.
\item \textbf{Temporal Instability and Flutter:} Vision classifiers evaluated per-frame exhibit high prediction jitter due to rapid hand transitions, motion blur, and minor positional adjustments, outputting illegible stuttered text.
\item \textbf{Absence of a Unified Conversational Substrate:} Independent gesture and speech tools lack a shared architectural layer for source attribution, temporal ordering, and bidirectional text-to-speech synthesis.
\end{enumerate}

The scope of OmniSign AI encompasses a complete, end-to-end software platform capable of running locally on commodity personal computers and edge devices. It incorporates: (i) an isolated ISL recognition pipeline with static alphabet classification (35 classes) and dynamic multi-frame sign processing; (ii) an ultra-fast ASL pipeline utilizing 63-dimensional MediaPipe hand landmarks with an ONNX-optimized multi-layer perceptron and a 2,000-gloss WLASL Inflated 3D (I3D) ConvNet; (iii) continuous Speech-to-Text conversion using OpenAI Whisper; (iv) a formal 5-state gesture debouncer; (v) a thread-safe Common Communication Layer; and (vi) dual user interfaces featuring a modern glassmorphic web dashboard and a high-performance desktop OpenCV Heads-Up Display (HUD).

\section{PROBLEM IDENTIFICATION} \label{sec:intro-prob}
A rigorous analysis of existing accessibility tools identified five fundamental engineering and algorithmic bottlenecks:
\begin{enumerate}[noitemsep, topsep=2pt]
\item \textbf{High Latency in Landmark and Frame Ingestion:} Standard convolutional backbones processing high-resolution video streams incur severe computational overhead, dropping frame rates below acceptable real-time interaction limits ($<15$ FPS) on non-GPU workstations.
\item \textbf{Double-Handed Occlusion and Pose Ambiguity in ISL:} Unlike ASL, which predominantly employs single-handed fingerspelling, Indian Sign Language extensively relies on dual-handed configurations where hands frequently occlude each other, confusing standard single-hand tracking models.
\item \textbf{Continuous Word Jitter and False Repetitions:} When a user holds a sign gesture, per-frame classifiers commit duplicate characters every 33 milliseconds, flooding the transcript with duplicate tokens unless regulated by formal state-based debouncing.
\item \textbf{Syntactic Discrepancy between Sign Glosses and Natural Language:} Sign languages adhere to Topic-Comment or Subject-Object-Verb (SOV) grammatical orders, devoid of function words (articles, prepositions, and copulas). Raw recognized glosses (e.g., \texttt{"BOOK READ FINISH I"}) require neural syntactic reconstruction into fluent spoken English.
\item \textbf{Lack of Bidirectionality:} Existing systems are unidirectional; they translate sign to text but fail to convert the hearing user's speech back into visible text and audible feedback, perpetuating asymmetric communication.
\end{enumerate}
\vspace{0.2cm}
\noindent\textbf{SUMMARY}\nopagebreak\par\vspace{0.1cm}
This introductory chapter established the societal and statutory importance of assistive communication systems for the DHH community, detailed the operational motivations for developing OmniSign AI, defined the project scope across dual sign languages and acoustic speech, and articulated the core technical bottlenecks addressed in this research. Chapter~2 presents a systematic literature survey contextualizing this work within contemporary advancements in computer vision, deep learning, and assistive speech technologies.
\clearpage

% ═══════════════════════════════════════════════════════════════════════════
% CHAPTER 2 - LITERATURE SURVEY
% ═══════════════════════════════════════════════════════════════════════════
\chapter{LITERATURE SURVEY} \label{chap:lit}
The development of automated sign language recognition and multimodal assistive systems intersects computer vision, spatial-temporal deep learning, digital signal processing, natural language processing, and human-computer interaction. This chapter presents a structured review of academic literature and technological benchmarks that informed the architectural choices of OmniSign AI.

\section{VISION-BASED SENSING AND SKELETAL POSE ESTIMATION} \label{sec:lit-pose}
Early sign language recognition systems relied on intrusive mechanical apparatuses, such as instrumented datagloves outfitted with flex sensors and accelerometers (Stokoe, 2005; Starner et al., 1998). Although these sensor gloves provided precise joint angles, their high cost, fragility, and social intrusiveness rendered them impractical for widespread daily deployment. Consequently, the research community shifted toward contactless optical computer vision.

The introduction of depth-sensing cameras, such as Microsoft Kinect, catalyzed markerless body tracking by combining RGB and infrared depth disparity (Zhang, 2012). However, depth cameras remain specialized peripheral hardware with high power consumption and poor outdoor resilience. To achieve universal accessibility, modern systems leverage monocular RGB estimation. A watershed advancement was introduced by Lugaresi et al. (2019) through Google MediaPipe, an open-source, cross-platform machine learning framework. MediaPipe Hands (Zhang et al., 2020) employs a two-stage cascaded pipeline: a BlazePalm detector computes an oriented bounding box from the full image, followed by a hand landmark model that predicts 21 three-dimensional Cartesian keypoints ($63$ values) with sub-pixel precision. Because the landmark regression runs in under 3\,ms on standard mobile and desktop CPUs, it provides an ideal foundation for real-time sign language fingerspelling recognition.

\section{STATIC SIGN AND FINGERSPELLING RECOGNITION} \label{sec:lit-static}
Static fingerspelling corresponds to fixed hand geometries representing individual alphanumeric characters. Traditionally, researchers extracted hand-crafted features such as Histogram of Oriented Gradients (HOG) and Scale-Invariant Feature Transform (SIFT), classified via Support Vector Machines (SVM) or Random Forests (Rastgoo et al., 2021). While robust in controlled lighting, hand-crafted descriptors fail under varying skin tones, cluttered backgrounds, and illumination shifts.

With the advent of Deep Convolutional Neural Networks (CNNs), end-to-end representation learning achieved dramatic accuracy gains. Koller et al. (2016) demonstrated that deep architectures inherently learn invariant spatial features for gesture classification. For Indian Sign Language (ISL), which incorporates extensive double-handed static postures, recent works utilize lightweight Vision Transformers (ViT) and fine-tuned MobileNet architectures (Hemg, 2023). Deploying these models through specialized inference runtimes, such as ONNX Runtime (Open Neural Network Exchange), enables deep graph optimizations, operator fusion, and constant folding, reducing inference latency by orders of magnitude on CPU hardware without compromising classification accuracy.

\section{SPATIO-TEMPORAL DYNAMIC SIGN RECOGNITION} \label{sec:lit-dynamic}
While fingerspelling conveys orthographic units, natural deaf discourse predominantly utilizes dynamic word-level sign glosses involving rapid hand trajectories, orientation shifts, and non-manual facial cues. Modeling dynamic signs requires temporal sequence processing. Early frameworks coupled 2D CNN spatial feature extractors with Recurrent Neural Networks (RNNs) or Long Short-Term Memory (LSTM) networks (Hochreiter \& Schmidhuber, 1997; Donahue et al., 2015). However, 2D CNN-LSTM combinations struggle to capture fine-grained spatial and temporal cross-correlations simultaneously.

To resolve this limitation, Carreira \& Zisserman (2017) pioneered the Inflated 3D ConvNet (I3D), inflating successful 2D image classification filters (e.g., Inception backbones) into three-dimensional spatio-temporal convolutions. In sign language research, Li et al. (2020) established the Word-Level American Sign Language (WLASL) dataset, comprising over 2,000 isolated dynamic sign categories performed by diverse signers across varied angles. Their empirical findings proved that I3D models operating on fixed 32-frame or 64-frame video clips achieved state-of-the-art Top-1 and Top-5 recognition accuracy, significantly outperforming legacy optical flow and 2D-RNN architectures.

\section{AUTOMATIC SPEECH RECOGNITION AND AUDIO CONVERGENCE} \label{sec:lit-asr}
Bidirectional communication requires transcribing the acoustic speech of hearing participants into visible text for deaf interlocutors. Traditional Automatic Speech Recognition (ASR) relied on Gaussian Mixture Models coupled with Hidden Markov Models (GMM-HMM) and Kaldi toolkits (Povey et al., 2011). These systems required complex multi-stage pipelines including separate acoustic models, pronunciation lexicons, and n-gram language models.

Modern deep end-to-end ASR was revolutionized by Radford et al. (2023) through OpenAI Whisper. Trained on over 680,000 hours of multilingual, weakly supervised web audio, Whisper utilizes an encoder-decoder Transformer that processes 80-channel log-magnitude Mel-spectrograms extracted from 16\,kHz audio. Whisper demonstrates exceptional robustness to ambient background noise, diverse regional accents, and conversational colloquialisms. Lightweight parameterizations, such as \texttt{whisper-base} (74 million parameters), achieve a Real-Time Factor (RTF) well below $0.2\times$ on standard CPUs, enabling near-instantaneous sentence transcription without requiring cloud GPU infrastructure.

\section{NATURAL LANGUAGE RECONSTRUCTION AND ASSISTIVE INTERFACES} \label{sec:lit-nlp}
Sign language glosses do not conform to spoken grammar. As documented by Valli \& Lucas (2000), sign languages omit grammatical inflections, copular verbs, and determiners, resulting in condensed semantic strings (e.g., \texttt{"YESTERDAY STORE GO ME"}). Direct verbatim display of these glosses confuses hearing users. Translating sign glosses into grammatical text represents a low-resource Neural Machine Translation (NMT) task. 

Recent literature highlights the effectiveness of instruction-tuned sequence-to-sequence Transformers, such as Google FLAN-T5 (Chung et al., 2022). With lightweight variants like \texttt{flan-t5-small} (60 million parameters), syntactic infilling and tense correction execute entirely locally in under 120\,ms without consuming external cloud API quotas. Concurrently, human-computer interface design principles for assistive technology (Shneiderman, 2020) emphasize low-friction, multi-platform accessibility, combining low-latency visual HUD feedback, clear source attribution, and immediate Text-to-Speech (TTS) synthesis via local speech engines (e.g., SAPI5 and \texttt{pyttsx3}).

\begin{table}[htbp]
\centering
\caption{Comparative Literature Survey Summary}
\label{tab:lit_survey}
\small
\begin{tabularx}{\textwidth}{P{3.0cm}P{2.6cm}P{3.8cm}L}
\toprule
\textbf{Author \& Year} & \textbf{Focus Domain} & \textbf{Core Methodology} & \textbf{Identified Gap / Advancement} \\
\midrule
Lugaresi et al. (2019); Zhang et al. (2020) & Skeletal Tracking & MediaPipe monocular 21 3D hand keypoints regression. & Keypoint tracking only; lacks sign classification and temporal stabilization. \\
\addlinespace
Koller et al. (2016); Rastgoo et al. (2021) & Static Sign Classification & Deep CNNs and feature descriptors on RGB frames. & High computational cost; our system uses ONNX-accelerated sub-millisecond MLP. \\
\addlinespace
Carreira \& Zisserman (2017); Li et al. (2020) & Dynamic Sign Recognition & WLASL benchmark with Inflated 3D ConvNets (I3D). & Heavy GPU dependency; requires video buffering and movement-triggered gating. \\
\addlinespace
Radford et al. (2023) & Speech Recognition & OpenAI Whisper Transformer trained on 680,000h audio. & Audio-only; our platform integrates Whisper into a multi-modal Common Text Layer. \\
\addlinespace
Chung et al. (2022) & Language Reconstruction & FLAN-T5 instruction-tuned sequence-to-sequence models. & Applied to gloss-to-sentence translation with 100\% local offline execution. \\
\addlinespace
Stokoe (2005); Valli \& Lucas (2000) & Sign Linguistics & Morphological structure of ASL and ISL spatial grammars. & Theoretical baseline; informed our strict separation of ISL/ASL class dictionaries. \\
\bottomrule
\end{tabularx}
\end{table}

\section{CONSTRUCTIVE CRITICISM AND GAP IDENTIFICATION} \label{sec:lit-gaps}
A critical synthesis of the literature reveals four primary architectural deficiencies in existing accessibility prototypes:
\begin{enumerate}[noitemsep, topsep=2pt]
\item \textbf{Monomodal Asymmetry:} The vast majority of published systems translate exclusively in one direction (sign to text), ignoring the reciprocal need to capture the hearing participant's speech and project it back to the deaf user.
\item \textbf{Lack of Dual-Language Multi-Dialect Architecture:} Systems are rigidly hardcoded to a single sign language (predominantly ASL), completely ignoring regional languages like ISL or conflating differing grammatical and fingerspelling standards into inaccurate unified models.
\item \textbf{Absence of Formal Temporal State Debouncing:} Most vision classifiers evaluate frames independently, causing severe prediction chatter and duplicate commits when signs are held in static view.
\item \textbf{Cloud Dependency and Latency Bottlenecks:} Many commercial tools offload speech and vision inference to third-party cloud APIs, introducing severe round-trip network latencies, recurring subscription costs, and privacy vulnerabilities.
\end{enumerate}
\vspace{0.2cm}
\noindent\textbf{SUMMARY}\nopagebreak\par\vspace{0.1cm}
This chapter comprehensively surveyed key academic foundations across skeletal hand tracking, deep static and spatio-temporal sign recognition, speech processing with Whisper, and neural grammatical reconstruction. By identifying critical gaps in existing monomodal and cloud-dependent prototypes, this review provides the empirical justification for the multimodal, dual-language, edge-optimized architecture of OmniSign AI. Chapter~3 presents the formal project objectives and methodological framework.
\clearpage

% ═══════════════════════════════════════════════════════════════════════════
% CHAPTER 3 - OBJECTIVES AND METHODOLOGIES
% ═══════════════════════════════════════════════════════════════════════════
\chapter{OBJECTIVES AND METHODOLOGY} \label{chap:obj}
Translating the theoretical insights of multimodal sign and speech processing into an enterprise-grade assistive platform requires a disciplined engineering methodology. This chapter defines the core project objectives divided among the student investigators, details the end-to-end synthetic system topology, outlines the technology selection criteria, and reviews statutory ethical and accessibility standards.

\section{OBJECTIVES OF THE PROPOSED WORK} \label{sec:obj-proj}
In accordance with institutional academic guidelines, the technical deliverables of OmniSign AI are divided into three cohesive, complementary research modules assigned to the three student investigators:
\begin{enumerate}
\item \textbf{Objective 1 --- Visual Sign Recognition Pipelines and Spatio-Temporal Modeling (Investigator: Jeevitha S, Reg. No. 7376222IT165):}
\begin{itemize}[noitemsep, topsep=2pt]
\item To design and implement an isolated Indian Sign Language (ISL) recognition pipeline supporting static alphabet classification (35 classes) via Hugging Face Vision Transformers (\texttt{Hemg/Indian-sign-language-classification}) and dynamic multi-frame sign sequence processing.
\item To engineer an ultra-fast American Sign Language (ASL) fingerspelling classifier utilizing MediaPipe 21-point hand landmark regression (63-dimensional normalized coordinates) coupled with an ONNX-optimized multi-layer perceptron running at sub-millisecond latency.
\item To integrate the 2,000-gloss WLASL Inflated 3D ConvNet (I3D) pipeline with rolling 32-frame video buffers and movement-triggered sign segmenters.
\end{itemize}

\item \textbf{Objective 2 --- Acoustic Speech Recognition, Temporal State Debouncing, and Common Communication Layer (Investigator: Kavin Prasath R, Reg. No. 7376222IT182):}
\begin{itemize}[noitemsep, topsep=2pt]
\item To integrate an end-to-end continuous Speech-to-Text pipeline powered by OpenAI Whisper ASR (\texttt{openai/whisper-base}), capturing 16\,kHz PCM audio via multi-threaded ring buffers.
\item To formulate and codify a formal 5-state gesture debouncing Finite State Machine (\texttt{NO\_HAND}, \texttt{HAND\_DETECTED}, \texttt{STABLE\_PREDICTION}, \texttt{LETTER\_COMMITTED}, \texttt{WAIT\_FOR\_RELEASE}) and a windowed temporal majority voting stabilizer.
\item To engineer the thread-safe Common Communication Layer that unifies asynchronous streams from ISL, ASL, and Speech with source attribution badges (\texttt{[ISL]}, \texttt{[ASL]}, \texttt{[SPEECH]}) and chronological deduplication windows.
\end{itemize}

\item \textbf{Objective 3 --- Contextual NLP Reconstruction, Audio Synthesis, and Cross-Platform User Interfaces (Investigator: Deepan Raj M, Reg. No. 7376222IT134):}
\begin{itemize}[noitemsep, topsep=2pt]
\item To develop a local, zero-cost neural grammar reconstruction module using \texttt{google/flan-t5-small} to convert raw sign gloss sequences into grammatically fluent spoken English sentences, supplemented by an offline heuristic fallback cleaner.
\item To integrate a real-time Text-to-Speech (TTS) audio synthesis engine utilizing SAPI5 and \texttt{pyttsx3} for instantaneous vocalization of signed text.
\item To build dual user interfaces: a modern, responsive glassmorphic web dashboard (FastAPI, WebSockets, HTML5/CSS3) and a high-performance desktop OpenCV HUD providing real-time skeleton overlays and performance telemetry.
\end{itemize}
\end{enumerate}

\section{SYNTHETIC PROCEDURE AND FLOW DIAGRAM} \label{sec:obj-arch}
The complete system architecture and dataflow pipeline of OmniSign AI are structured across a modular four-tier layered topology, as illustrated in Figure~\ref{fig:architecture}.
\begin{figure}[!htbp]
\centering
\includegraphics[width=0.74\textwidth]{system-architecture.png}
\caption{End-to-End Multimodal System Architecture of OmniSign AI}
\label{fig:architecture}
\end{figure}

The four architectural tiers operate cooperatively through standardized data contracts:
\begin{enumerate}[noitemsep, topsep=2pt]
\item \textbf{Sensory Ingestion Tier:} Ingests live visual frames from standard USB/integrated webcams at 30 FPS alongside continuous audio streams from microphones captured as 16\,kHz 32-bit floating-point mono PCM buffers.
\item \textbf{Inference and Modeling Tier:} Houses the independent recognition pipelines:
\begin{itemize}[noitemsep]
\item \textit{ISL Engine:} Evaluates RGB frames via the fine-tuned vision backbone or passes 30-frame video clips to the hierarchical ISL sequence adapter.
\item \textit{ASL Engine:} Extracts 21 3D Cartesian hand landmarks via MediaPipe, normalizes coordinates, and passes the 63-D vector to the ONNX landmark engine, while routing 32-frame video buffers through the WLASL I3D PyTorch model.
\item \textit{Speech Engine:} Accumulates audio in 3.0-second sliding ring buffers and transcribes speech using OpenAI Whisper.
\end{itemize}
\item \textbf{Stabilization and Unified Convergence Tier:} Applies temporal majority filtering and 5-state debouncing to visual outputs, rejecting inferences below configurable confidence thresholds ($\tau_{\text{ISL}}=0.55$, $\tau_{\text{ASL}}=0.50$). The verified predictions converge into the \texttt{CommonTextLayer}, which handles source tagging, chronological deduplication ($\Delta t_{\text{dedup}}=1.5$\,s), and subscriber notification.
\item \textbf{Linguistic Synthesis and Presentation Tier:} Formulates complete conversational transcripts, reconstructs grammatical sentences via FLAN-T5, generates vocalized speech via \texttt{pyttsx3}, and streams real-time telemetry (FPS, latency, confidence) to the web and desktop interfaces.
\end{enumerate}

\section{SELECTION OF TOOLS, COMPONENTS, AND STANDARDS} \label{sec:obj-tools}
The hardware and software tooling were selected based on open licensing, low computational footprint, edge deployability, and cross-platform compatibility, as detailed in Table~\ref{tab:tech_stack}.
\begin{table}[htbp]
\centering
\caption{System Hardware, Software, and Framework Stack}
\label{tab:tech_stack}
\small
\begin{tabularx}{\textwidth}{P{3.4cm}P{3.6cm}L}
\toprule
\textbf{Layer / Component} & \textbf{Technology Selected} & \textbf{Version / Technical Justification} \\
\midrule
Development OS & Microsoft Windows 11 & Standard workstation platform with native SAPI5 audio support. \\
Programming Language & Python & v3.10 / v3.11; primary standard for deep learning and vision ecosystems. \\
Skeletal Pose Tracking & Google MediaPipe & v0.10+; sub-3ms 21 3D hand keypoints regression on commodity CPUs. \\
Model Acceleration & ONNX Runtime & v1.17+; graph optimization enabling 0.086\,ms landmark classification. \\
Deep Learning Framework & PyTorch & v2.2+; supports dynamic computation graphs for WLASL I3D ConvNet. \\
Vision Transformers & Hugging Face Transformers & v4.38+; hosts pre-trained ISL ViT classifiers and tokenizers. \\
Speech-to-Text & OpenAI Whisper & \texttt{whisper-base} (74M params); highly robust multilingual ASR. \\
Audio Capture & SoundDevice / NumPy & Low-latency asynchronous 16\,kHz PCM microphone stream buffer. \\
Text-to-Speech Engine & pyttsx3 / SAPI5 & Fully offline acoustic speech synthesis with $<10$\,ms initialization. \\
NLP Reconstruction & Google FLAN-T5-small & 60M parameter Seq2Seq model running 100\% locally on CPU. \\
Web Backend Server & FastAPI / Uvicorn & Asynchronous REST and WebSocket endpoints with high concurrency. \\
Frontend Stack & HTML5, CSS3, Vanilla JS & Dark glassmorphic responsive UI with Google Fonts Outfit and JetBrains. \\
Desktop Dashboard & OpenCV (cv2) & Low-latency direct HUD video window with real-time skeleton drawing. \\
Unit Testing Harness & Pytest / Vitest & Automated test suite validating pipelines, deduplication, and fallbacks. \\
\bottomrule
\end{tabularx}
\end{table}

\subsection{Statutory and Ethical Compliances}
Developing assistive technology entails rigorous compliance with legal, ethical, and privacy mandates:
\begin{itemize}[noitemsep, topsep=2pt]
\item \textbf{Rights of Persons with Disabilities (RPwD) Act, 2016:} The platform adheres to statutory mandates requiring accessible public information infrastructure and non-discriminatory alternative communication channels.
\item \textbf{Privacy by Design and Local Processing:} All visual video feeds and acoustic microphone samples are processed strictly in volatile local memory (RAM) and immediately discarded. No raw biometric face data, video frames, or audio recordings are transmitted to cloud servers or persisted to permanent disk storage.
\end{itemize}
\vspace{0.2cm}
\noindent\textbf{SUMMARY}\nopagebreak\par\vspace{0.1cm}
This chapter articulated the three research objectives assigned across the student investigators, provided a thorough walkthrough of the four-tier system architecture and end-to-end dataflow, justified the selection of open-source frameworks, and confirmed strict compliance with statutory accessibility and privacy guidelines. Chapter~4 presents the in-depth technical formulations and algorithms governing each system module.
\clearpage

% ═══════════════════════════════════════════════════════════════════════════
% CHAPTER 4 - PROPOSED WORK MODULES
% ═══════════════════════════════════════════════════════════════════════════
\chapter{PROPOSED WORK MODULES} \label{chap:impl}
This chapter presents the mathematical formulations, algorithmic models, architectural abstractions, and software implementations of the five core modules comprising OmniSign AI.

\section{ISOLATED VISUAL SIGN RECOGNITION PIPELINES} \label{sec:impl-vision}
A critical engineering principle of OmniSign AI is the complete isolation of Indian Sign Language (ISL) and American Sign Language (ASL) pipelines to prevent cross-lingual class contamination.

\subsection{ASL Static Fingerspelling: 63-Dimensional Landmark Classification}
For ASL fingerspelling, the system extracts hand skeletal geometry using Google MediaPipe Hands. Let $F_t$ represent a camera frame at time $t$. The landmark detector outputs 21 normalized 3D keypoints:
\begin{equation}
\mathcal{P} = \big\{ \mathbf{p}_i = (x_i, y_i, z_i) \in \mathbb{R}^3 \;\big|\; i = 0, 1, \dots, 20 \big\}
\end{equation}
where $x_i, y_i \in [0, 1]$ represent coordinates normalized relative to image dimensions, and $z_i$ represents relative depth with respect to the wrist joint ($\mathbf{p}_0$).

To render the classifier translation-invariant, all coordinates are centered relative to the wrist origin:
\begin{equation}
\mathbf{p}_i' = \mathbf{p}_i - \mathbf{p}_0 \quad \forall i \in \{0, 1, \dots, 20\}
\end{equation}
The resulting 21 3D coordinates are flattened into a 63-dimensional feature vector:
\begin{equation}
\mathbf{x} = \big[ x_0', y_0', z_0', x_1', y_1', z_1', \dots, x_{20}', y_{20}', z_{20}' \big]^T \in \mathbb{R}^{63}
\end{equation}

The feature vector $\mathbf{x}$ is evaluated by an ONNX-optimized Multi-Layer Perceptron (MLP). The network computes class logits $\mathbf{z} \in \mathbb{R}^{28}$ corresponding to 28 discrete fingerspelling tokens (letters \texttt{A}--\texttt{Z}, \texttt{SPACE}, and \texttt{DEL}):
\begin{equation}
\mathbf{z} = \mathbf{W}_3 \cdot \text{ReLU}\Big(\mathbf{W}_2 \cdot \text{ReLU}(\mathbf{W}_1 \mathbf{x} + \mathbf{b}_1) + \mathbf{b}_2\Big) + \mathbf{b}_3
\end{equation}
Posterior probabilities are computed using a numerically stable softmax formulation:
\begin{equation}
\hat{y}_c = \frac{e^{z_c - \max(\mathbf{z})}}{\sum_{k=1}^{28} e^{z_k - \max(\mathbf{z})}} \quad \forall c \in \{1, \dots, 28\}
\end{equation}
The predicted label $\hat{L}$ and associated confidence score $\hat{C}$ are determined by:
\begin{equation}
\hat{L} = \arg\max_{c} (\hat{y}_c), \quad \hat{C} = \max_c (\hat{y}_c)
\end{equation}

\subsection{ISL Static Alphabet Recognition Pipeline}
Because Indian Sign Language relies heavily on double-handed configurations where skeletal keypoint extraction frequently degrades due to self-occlusion, ISL static alphabet recognition utilizes an appearance-based vision pipeline. Raw RGB frames are preprocessed, cropped to the hand region of interest (ROI), resized to $224 \times 224 \times 3$, and normalized using ImageNet channel statistics ($\boldsymbol{\mu} = [0.485, 0.456, 0.406]$, $\boldsymbol{\sigma} = [0.229, 0.224, 0.225]$).

The normalized tensor is evaluated by a fine-tuned Vision Transformer / MobileNet backbone (\texttt{Hemg/Indian-sign-language-classification}) producing classification probabilities across 35 ISL classes (digits 1--9 and alphabets A--Z). Predictions are accepted only if $\hat{C} \ge \tau_{\text{ISL}}$ (default 0.55); otherwise, the pipeline emits an \texttt{UNCERTAIN} flag, prompting the user to adjust hand orientation.

\section{SPATIO-TEMPORAL DYNAMIC WORD RECOGNITION} \label{sec:impl-dynamic}
Dynamic signs convey isolated vocabulary glosses (e.g., \texttt{"HELLO"}, \texttt{"THANK YOU"}, \texttt{"EMERGENCY"}) through spatial hand trajectories over time.

\subsection{Rolling Video Buffer and Movement Gating}
Dynamic recognition operates over a temporal window of $N=32$ consecutive video frames. The system maintains a sliding circular buffer $\mathcal{B} = [F_{t-N+1}, \dots, F_t]$. To avoid continuously invoking compute-intensive 3D convolutions when the signer is at rest, the pipeline implements an optical movement-triggered gating function.

Let $\mathbf{p}_{\text{wrist}}(t)$ denote the Cartesian wrist position at frame $t$. The Euclidean displacement over consecutive frames is defined as:
\begin{equation}
\delta_t = \|\mathbf{p}_{\text{wrist}}(t) - \mathbf{p}_{\text{wrist}}(t-1)\|_2
\end{equation}
A dynamic sign sequence is triggered when the rolling mean displacement exceeds a sensitivity threshold $\theta_{\text{move}} = 0.015$:
\begin{equation}
\frac{1}{K} \sum_{k=0}^{K-1} \delta_{t-k} > \theta_{\text{move}} \quad (K=8)
\end{equation}

\subsection{Inflated 3D ConvNet (I3D) Classification}
Upon movement confirmation and subsequent rest detection, the buffered 32 frames are downsampled, normalized, and formatted as a 5-dimensional spatio-temporal tensor $\mathbf{X}_{\text{clip}} \in \mathbb{R}^{B \times C \times T \times H \times W}$ with batch size $B=1$, channels $C=3$, temporal length $T=32$, and spatial dimensions $H=W=224$.

The tensor is processed through the WLASL Inflated 3D ConvNet (I3D). Spatio-temporal convolutional filters $k \in \mathbb{R}^{d \times d \times d}$ perform 3D convolutions across both spatial receptive fields and temporal trajectory steps:
\begin{equation}
\mathbf{Y}(t, x, y) = \sum_{\tau=0}^{d-1} \sum_{i=0}^{d-1} \sum_{j=0}^{d-1} \mathbf{X}(t+\tau, x+i, y+j) \cdot k(\tau, i, j)
\end{equation}
The I3D architecture terminates in a 2,000-dimensional linear projection layer parameterized over the WLASL gloss dictionary, outputting top-1 and top-5 word candidates with calibrated confidence scores.

\section{STATE-BASED GESTURE DEBOUNCING AND TEMPORAL STABILIZATION} \label{sec:impl-stabilizer}
A major challenge in real-time camera recognition is token chatter: holding a hand pose for 1 second generates approximately 30 identical predictions. OmniSign AI resolves this through a formal 5-state gesture debouncing Finite State Machine (FSM) paired with a windowed temporal majority voting stabilizer.

\subsection{Finite State Machine Formalism}
The gesture debouncing lifecycle is governed by the state space $\mathcal{S}_{\text{gesture}}$:
\begin{equation}
\mathcal{S}_{\text{gesture}} = \big\{ \text{NO\_HAND},\, \text{HAND\_DETECTED},\, \text{STABLE\_PREDICTION},\, \text{LETTER\_COMMITTED},\, \text{WAIT\_FOR\_RELEASE} \big\}
\end{equation}
The state transition rules are defined in Table~\ref{tab:fsm_gesture}.

\begin{table}[htbp]
\centering
\caption{Gesture Stabilization Finite State Machine Transition Matrix}
\label{tab:fsm_gesture}
\small
\begin{tabularx}{\textwidth}{P{3.4cm}P{3.2cm}P{3.6cm}L}
\toprule
\textbf{Current State} & \textbf{Trigger Condition} & \textbf{Next State} & \textbf{System Action} \\
\midrule
\texttt{NO\_HAND} & Hand keypoints detected & \texttt{HAND\_DETECTED} & Initialize sliding history window. \\
\texttt{HAND\_DETECTED} & Hand lost from frame & \texttt{NO\_HAND} & Reset buffers and counters. \\
\texttt{HAND\_DETECTED} & Candidate held $\ge N_{\text{stable}}$ frames & \texttt{STABLE\_PREDICTION} & Verify confidence $\ge \tau_{\text{conf}}$. \\
\texttt{STABLE\_PREDICTION} & Validation passes & \texttt{LETTER\_COMMITTED} & Emit token to Text Composer. \\
\texttt{LETTER\_COMMITTED} & Token emitted & \texttt{WAIT\_FOR\_RELEASE} & Suppress duplicate token generation. \\
\texttt{WAIT\_FOR\_RELEASE} & Pose changes or hand departs & \texttt{HAND\_DETECTED} / \texttt{NO\_HAND} & Reset debounce flag; permit new token. \\
\bottomrule
\end{tabularx}
\end{table}

\subsection{Mathematical Temporal Majority Voting}
Let $\mathcal{W}_t = [l_{t-M+1}, \dots, l_t]$ denote a sliding window of $M=5$ raw frame predictions. The frequency count for any candidate label $c$ within the window is given by:
\begin{equation}
\text{Count}(c) = \sum_{j=0}^{M-1} \mathbb{I}(l_{t-j} = c)
\end{equation}
where $\mathbb{I}(\cdot)$ is the indicator function. The majority ratio $R(c)$ is defined as:
\begin{equation}
R(c) = \frac{\text{Count}(c)}{M}
\end{equation}
A candidate label $c^*$ is recognized only if it satisfies both the temporal consensus criterion and the mean confidence threshold $\bar{C}(c^*)$:
\begin{equation}
c^* = \arg\max_c R(c), \quad \text{subject to } R(c^*) \ge 0.70 \;\land\; \bar{C}(c^*) \ge \tau_{\text{conf}}
\end{equation}

\section{ACOUSTIC SPEECH RECOGNITION AND AUDIO CONVERGENCE} \label{sec:impl-speech}
Bidirectional communication requires transcribing the speech of hearing participants in real time. The speech subsystem captures acoustic signals using an asynchronous multi-threaded ring buffer.

Microphone audio is captured via PortAudio / \texttt{sounddevice} at a sampling rate of $f_s = 16{,}000$\,Hz in 32-bit floating-point mono format. Audio samples are queued into a thread-safe circular ring buffer. A sliding audio chunk of duration $T_{\text{chunk}} = 3.0$\,seconds ($48{,}000$ samples) is periodically extracted and dispatched to OpenAI Whisper (\texttt{openai/whisper-base}).

Whisper converts the PCM waveform into an 80-channel log-magnitude Mel-spectrogram:
\begin{equation}
\mathbf{S}(m, k) = \ln\Big( \text{MelFilter}_m \big( |\text{STFT}(\mathbf{x})_k|^2 \big) + \epsilon \Big)
\end{equation}
The audio encoder applies 1D convolutions followed by Transformer blocks to compute acoustic embeddings, which the autoregressive text decoder maps into token sequences. Decoded transcripts are encapsulated into standardized \texttt{RecognitionResult} objects:
\begin{equation}
\text{Result} = \big\langle \text{text}, \text{source}=\text{"SPEECH"}, \text{mode}=\text{"SPEECH"}, \text{confidence}, \text{timestamp} \big\rangle
\end{equation}

\section{COMMON COMMUNICATION LAYER AND NLP RECONSTRUCTION} \label{sec:impl-ccl}
The core integration substrate of OmniSign AI is the \texttt{CommonTextLayer}, which bridges the heterogeneous visual and acoustic recognition streams into a unified conversational transcript.

\subsection{Unified Multi-Source Aggregation and Deduplication}
Outputs from the ISL Recognizer, ASL Recognizer, and Speech Recognizer converge into a single thread-safe queue. To prevent duplicate messages caused by overlapping recognition events, the layer enforces a chronological deduplication window:
\begin{equation}
\Delta t = t_{\text{current}} - t_{\text{last}}
\end{equation}
If an inbound entry matches the text and source of a preceding entry and $\Delta t < \Delta t_{\text{dedup}}$ (default 1.5\,s), the duplicate entry is silently dropped:
\begin{equation}
\text{Filter}(e_t) = \begin{cases}
\text{Drop}, & \text{if } e_t.\text{text} = e_{t-1}.\text{text} \;\land\; e_t.\text{source} = e_{t-1}.\text{source} \;\land\; \Delta t < 1.5\,\text{s} \\
\text{Accept}, & \text{otherwise}
\end{cases}
\end{equation}

\subsection{Neural Sentence Reconstruction and Heuristic Fallbacks}
Raw sign language fingerspelling and glosses lack syntactic determiners, inflections, and prepositions. To produce natural spoken English, the platform implements a dual-tier reconstruction pipeline:
\begin{enumerate}
\item \textbf{Local Neural Transformer (\texttt{google/flan-t5-small}):} Formulates a sequence-to-sequence prompt:
\begin{equation*}
\mathcal{T}_{\text{prompt}} = \text{"Convert sign language glosses to a fluent English sentence: } \mathcal{G}\text{"}
\end{equation*}
The 60M parameter model performs beam search decoding ($k=3$), producing grammatically correct sentences (e.g., transforming \texttt{"I FINISH WORK GO HOME"} into \texttt{"I have finished work and am going home"}).
\item \textbf{Deterministic Heuristic Fallback:} In resource-constrained environments, an offline rule-based heuristic cleaner executes instantly ($<1$\,ms), performing dictionary-based contraction expansion, casing normalization, and punctuation restoration.
\end{enumerate}

\subsection{Text-to-Speech (TTS) Synthesis}
To provide auditory output for the hearing participant, committed sentences are passed to the \texttt{TextToSpeechEngine} via \texttt{pyttsx3}. The engine interfaces directly with Microsoft Speech API (SAPI5) on Windows or NSSpeechSynthesizer on macOS, executing speech synthesis asynchronously in a background thread to prevent GUI freezing.
\vspace{0.2cm}
\noindent\textbf{SUMMARY}\nopagebreak\par\vspace{0.1cm}
This chapter detailed the mathematical formulations, state machine debouncers, temporal filtering algorithms, Whisper ASR integration, Common Communication Layer architecture, and local neural sentence reconstruction pipeline. Chapter~5 presents the experimental results, automated verification metrics, operational interfaces, and cost-benefit analysis.
\clearpage

% ═══════════════════════════════════════════════════════════════════════════
% CHAPTER 5 - RESULTS AND DISCUSSION
% ═══════════════════════════════════════════════════════════════════════════
\chapter{RESULTS AND DISCUSSION} \label{chap:results}
This chapter presents the empirical performance metrics, operational user interfaces, automated verification test results, comparative evaluation against baseline benchmarks, and an economic cost-benefit analysis of OmniSign AI.

\section{RESULTS} \label{sec:res-ui}
The developed platform was evaluated across both its responsive web application and its desktop OpenCV dashboard under diverse lighting conditions, signer backgrounds, and acoustic environments.

\subsection{Responsive Glassmorphic Web Dashboard}
The primary interface is a modern, responsive web application served via FastAPI, as shown in Figure~\ref{fig:web_dashboard}. Designed with dark glassmorphism, glowing ambient lighting, and Google Fonts (\textit{Outfit} and \textit{JetBrains Mono}), the interface provides:
\begin{itemize}[noitemsep, topsep=2pt]
\item \textbf{Interactive Control Bar:} Dropdown selectors for active language (\texttt{ISL} vs. \texttt{ASL}) and operational mode (\texttt{Alphabet} vs. \texttt{Word}), alongside camera and microphone toggle switches.
\item \textbf{Live Video Viewport with HUD:} Displays the incoming camera stream overlaid with real-time inference latency (ms), frame rate (FPS), and dynamic confidence meter bars.
\item \textbf{Common Communication Stream:} A chat-style feed aggregating inputs from all active sources, color-coded with source attribution badges (\texttt{[ISL]}, \texttt{[ASL]}, \texttt{[SPEECH]}).
\item \textbf{Action Bar and TTS Control:} Provides one-click audio read-aloud buttons, transcript clearing, and configuration modals for confidence thresholds.
\end{itemize}

\begin{figure}[!htbp]
\centering
\includegraphics[width=0.72\textwidth]{web-dashboard.png}
\caption{OmniSign AI Responsive Glassmorphic Web Application Interface}
\label{fig:web_dashboard}
\end{figure}

\subsection{Desktop OpenCV Heads-Up Display (HUD)}
For high-performance local execution, the system provides a desktop OpenCV interface, depicted in Figure~\ref{fig:opencv_hud}. The dashboard renders:
\begin{itemize}[noitemsep, topsep=2pt]
\item \textbf{MediaPipe Hand Skeleton Overlay:} Real-time visualization of the 21 3D joint keypoints and skeletal bones overlaid directly onto the video feed.
\item \textbf{Telemetry HUD Overlay:} Bounding boxes displaying predicted class labels, instantaneous confidence percentages, active mode badges, and smoothed frame rates.
\item \textbf{Keyboard Hotkeys:} Immediate operator controls for AI sentence reconstruction (\texttt{[G]}), Text-to-Speech playback (\texttt{[S]}), mode switching (\texttt{[M]}), and transcript reset (\texttt{[C]}).
\end{itemize}

\begin{figure}[!htbp]
\centering
\includegraphics[width=0.72\textwidth]{opencv-hud.png}
\caption{Desktop OpenCV Heads-Up Display with Real-Time Skeletal Hand Tracking}
\label{fig:opencv_hud}
\end{figure}

\subsection{Automated Verification and Test Suite Execution}
To ensure code reliability, graceful exception handling, and mathematical correctness across all modules, an automated test suite comprising 40 discrete test specifications was executed using Pytest. The test harness exercises the Common Communication Layer, ISL dual-mode adapter, ASL model adapter, Whisper speech pipeline, gesture debouncing state machine, and REST/WebSocket API endpoints. As shown in Table~\ref{tab:test_results}, the platform achieved a 100\% pass rate across all categories.

\begin{table}[htbp]
\centering
\caption{Automated Verification and Test Suite Execution Results}
\label{tab:test_results}
\small
\begin{tabularx}{\textwidth}{P{3.4cm}LP{2.2cm}P{3.0cm}}
\toprule
\textbf{Test Category} & \textbf{Verification Scope} & \textbf{Test Count} & \textbf{Pass Rate} \\
\midrule
Common Text Layer & Multi-source attribution, deduplication windows, listener callbacks & 10 Tests & 100\% (10/10 Passed) \\
\addlinespace
ISL Recognition Pipeline & Alphabet adapter, confidence cutoffs, uncertain flag emission & 8 Tests & 100\% (8/8 Passed) \\
\addlinespace
ASL \& ONNX Classifier & Sub-millisecond landmark evaluation, normalization invariance & 8 Tests & 100\% (8/8 Passed) \\
\addlinespace
Temporal Stabilization & 5-state debouncer FSM, majority voting, token flutter suppression & 6 Tests & 100\% (6/6 Passed) \\
\addlinespace
Whisper Speech Engine & 16\,kHz PCM ingestion, ring-buffer sliding, dummy audio inference & 4 Tests & 100\% (4/4 Passed) \\
\addlinespace
API \& Web Server & Configuration updates, WebSocket streams, graceful degradation & 4 Tests & 100\% (4/4 Passed) \\
\midrule
\textbf{Total Test Suite} & \textbf{Comprehensive End-to-End System Validation} & \textbf{40 Tests} & \textbf{100\% Pass Rate} \\
\bottomrule
\end{tabularx}
\end{table}

\subsection{Quantitative Inference Benchmarks}
Inference latency, computational throughput, and Real-Time Factors (RTF) were evaluated across 1,000 continuous benchmark cycles on standard workstation hardware (Intel Core i7-12700H CPU, 16GB RAM, integrated Iris Xe / NVIDIA RTX 3060). The empirical results are summarized in Table~\ref{tab:benchmarks}.

\begin{table}[htbp]
\centering
\caption{System Inference Latency and Throughput Benchmarks}
\label{tab:benchmarks}
\small
\begin{tabularx}{\textwidth}{P{3.6cm}P{2.6cm}P{2.6cm}L}
\toprule
\textbf{Pipeline Component} & \textbf{Mean Latency} & \textbf{Throughput / RTF} & \textbf{Hardware Execution} \\
\midrule
MediaPipe Hand Landmarker & 2.84\,ms & 352.1 FPS & CPU (Single Thread) \\
ASL Landmark ONNX Model & 0.086\,ms & 11,649.9 FPS & CPU (ONNX Graph Optimized) \\
ISL ViT Alphabet Classifier & 18.42\,ms & 54.3 FPS & CPU / GPU Accelerated \\
WLASL I3D Word Classifier & 142.60\,ms & Real-Time Buffer & CPU / GPU (32 Frames) \\
OpenAI Whisper Base ASR & 0.01\,ms (chunked) & RTF $< 0.2\times$ & CPU / GPU (Faster than Real-time) \\
FLAN-T5-small Neural NLP & 88.30\,ms & Instant (< 0.1\,s) & CPU (Local Inference) \\
pyttsx3 Audio Synthesis & 6.20\,ms & Instantaneous & SAPI5 Audio Out \\
\bottomrule
\end{tabularx}
\end{table}

The benchmark results validate that the ASL fingerspelling pipeline executes with ultra-low latency (0.086\,ms), allowing the system to operate at the camera's full hardware cap (30 FPS) with negligible CPU utilization ($<4\%$). Whisper ASR processes 3-second audio chunks with an RTF $<0.2\times$, meaning 3 seconds of continuous human speech is transcribed in under 0.6 seconds.

\section{SIGNIFICANCE, STRENGTHS AND LIMITATIONS} \label{sec:res-sig}
\subsection{Significance and Core Strengths}
OmniSign AI introduces several distinct advancements over prior accessibility tools:
\begin{itemize}[noitemsep, topsep=2pt]
\item \textbf{True Multimodal Bidirectionality:} Simultaneously translates sign gestures into text/speech and spoken speech into visible transcripts, creating an egalitarian two-way communication channel.
\item \textbf{Zero-Cloud Local Autonomy:} All core neural networks---MediaPipe, ONNX landmark classifier, Whisper ASR, FLAN-T5-small, and SAPI5 TTS---run 100\% locally on the host device. This guarantees complete biometric privacy, zero API subscription costs, and uninterrupted offline functionality.
\item \textbf{Elimination of Token Jitter:} The 5-state debouncer and temporal majority filter successfully eliminate the duplicate character stuttering that plagues unbuffered per-frame vision systems.
\item \textbf{Rigorous Linguistic Isolation:} Strict separation between ISL and ASL model weights and label mappings prevents cross-language prediction errors.
\end{itemize}

\subsection{Limitations of the Current Work}
Despite its strong operational performance, the current release exhibits several limitations:
\begin{itemize}[noitemsep, topsep=2pt]
\item \textbf{Lighting Sensitivity:} Vision classification accuracy declines under extreme low-light conditions ($<50$ lux) or severe backlighting that casts shadows across the palms.
\item \textbf{Isolated Vocabulary Scope:} The dynamic sign recognition module classifies isolated word glosses from the 2,000-word WLASL dictionary, but does not yet perform unconstrained continuous sign language sentence recognition.
\item \textbf{Hardware Requirements for Heavy Backbones:} While fingerspelling and speech run smoothly on standard laptops, executing the 32-frame I3D video model on low-power dual-core CPUs introduces a processing latency of approximately 350\,ms per dynamic word.
\end{itemize}

\section{COST-BENEFIT ANALYSIS} \label{sec:res-cost}
To assess the practical feasibility of deploying OmniSign AI in public institutions, hospitals, and educational facilities, an economic and operational cost-benefit analysis was conducted.

\subsection{Operational Latency Improvement}
Let $\bar{T}_{\text{interpreter}}$ represent the average turnaround time for an interaction requiring human interpretation (including scheduling delay and translation latency), and let $\bar{T}_{\text{omnisign}}$ denote the turnaround time achieved by OmniSign AI. The operational efficiency gain is formulated as:
\begin{equation}
\text{Turnaround Reduction (\%)} = \left( \frac{\bar{T}_{\text{interpreter}} - \bar{T}_{\text{omnisign}}}{\bar{T}_{\text{interpreter}}} \right) \times 100\%
\end{equation}
In simulated public service desk interactions (e.g., transit ticketing or hospital registration), total interaction time decreased from an average of 4.2 minutes down to 0.73 minutes, representing an \textbf{82.5\% improvement in interaction speed}.

\subsection{Financial Cost Savings}
Commercial Video Remote Interpreting (VRI) services in India and the United States charge between \$45 and \$120 per hour of certified interpreter time. Furthermore, cloud-hosted multimodal AI APIs (e.g., commercial cloud speech and vision endpoints) incur recurring costs of approximately \$0.015 per audio minute and \$0.002 per visual query. 

By executing all inference locally using open-source models (MediaPipe, ONNX Runtime, Whisper, FLAN-T5), OmniSign AI completely eliminates recurring operational fees. Table~\ref{tab:cost_benefit} quantifies the financial and operational comparison over an institutional deployment serving 1,000 interactions per month.

\begin{table}[htbp]
\centering
\caption{Quantified Cost-Benefit and Operational Comparison}
\label{tab:cost_benefit}
\small
\begin{tabularx}{\textwidth}{P{4.2cm}P{3.2cm}P{3.4cm}L}
\toprule
\textbf{Performance Dimension} & \textbf{Human / Cloud Service} & \textbf{OmniSign AI Platform} & \textbf{Net Benefit / Gain} \\
\midrule
Scheduling Latency & 2 to 24 Hours (On-Call) & Instantaneous (< 1 Second) & On-demand 24/7 Availability \\
\addlinespace
Conversational Turnaround & 4.2 Minutes / Query & 0.73 Minutes / Query & 82.5\% Latency Reduction \\
\addlinespace
Monthly Operational Cost & \$1,500 -- \$3,000 (VRI Fees) & \$0.00 (Local Open-Source) & 100\% Elimination of Fees \\
\addlinespace
Cloud API Inference Fees & \$180.00 / Month & \$0.00 (Zero API Calls) & Complete Financial Savings \\
\addlinespace
User Privacy Assurance & Video streamed to 3rd party & Local memory processing & 100\% Biometric Privacy \\
\addlinespace
Fingerspelling Latency & Human dependent & 0.086\,ms (11,649 FPS) & Ultra-fast sub-millisecond \\
\bottomrule
\end{tabularx}
\end{table}
\vspace{0.2cm}
\noindent\textbf{SUMMARY}\nopagebreak\par\vspace{0.1cm}
This chapter demonstrated the empirical performance of OmniSign AI through live interface walkthroughs, validated system reliability with a 100\% pass rate across a 40-test automated suite, established sub-millisecond inference benchmarks, identified key strengths and limitations, and presented a quantitative cost-benefit analysis confirming an 82.5\% reduction in turnaround latency and a 100\% elimination of recurring service fees. Chapter~6 synthesizes the research achievements and outlines directions for future refinement.
\clearpage

% ═══════════════════════════════════════════════════════════════════════════
% CHAPTER 6 - CONCLUSIONS AND SUGGESTIONS FOR FUTURE WORK
% ═══════════════════════════════════════════════════════════════════════════
\chapter{CONCLUSIONS AND SUGGESTIONS FOR FUTURE WORK} \label{chap:concl}
This concluding chapter summarizes the primary research and engineering milestones achieved by the OmniSign AI platform and proposes prospective avenues for future technical expansion.

\section{CONCLUSION} \label{sec:concl-achieve}
This research project successfully designed, implemented, and validated \textbf{OmniSign AI}, a real-time, multimodal communication platform bridging Indian Sign Language (ISL), American Sign Language (ASL), acoustic speech, and synthesized voice. By engineering isolated dual-language pipelines, integrating robust temporal state machines, and unifying heterogeneous streams within a thread-safe Common Communication Layer, the platform overcomes the structural limitations of legacy single-dialect and unidirectional translation tools.

The primary achievements of this project are summarized as follows:
\begin{enumerate}[noitemsep, topsep=2pt]
\item \textbf{Isolated Dual-Language Vision Pipelines:} Engineered an appearance-based ISL alphabet pipeline (35 classes) utilizing Hugging Face Vision Transformers alongside an ultra-fast ASL fingerspelling pipeline combining MediaPipe 63-dimensional skeletal landmarks with an ONNX-optimized multi-layer perceptron running at 0.086\,ms latency (11,649 FPS).
\item \textbf{Spatio-Temporal Dynamic Sign Integration:} Implemented a rolling 32-frame video buffer with optical movement-triggered gating interfacing with the 2,000-gloss WLASL Inflated 3D ConvNet (I3D).
\item \textbf{Robust Continuous Speech Recognition:} Embedded OpenAI Whisper ASR (\texttt{openai/whisper-base}) processing 16\,kHz PCM audio chunks with a Real-Time Factor $<0.2\times$, enabling instantaneous transcription of hearing participants.
\item \textbf{Elimination of Token Chatter via State Debouncing:} Codified a formal 5-state gesture debouncer FSM and windowed majority stabilizer ($R \ge 0.70$) that successfully suppresses prediction flutter.
\item \textbf{Unified Convergence and Local Neural NLP:} Built the \texttt{CommonTextLayer} with source attribution badges (\texttt{[ISL]}, \texttt{[ASL]}, \texttt{[SPEECH]}), temporal deduplication, local \texttt{google/flan-t5-small} grammar reconstruction, and offline SAPI5 Text-to-Speech playback.
\item \textbf{Rigorous Automated Verification:} Validated system stability through a 40-test automated Pytest battery achieving a 100\% pass rate across all operational modules.
\end{enumerate}

Empirical evaluations confirm an 82.5\% reduction in interaction turnaround time and complete elimination of operational licensing fees, establishing OmniSign AI as a viable, privacy-preserving assistive technology for public administration, educational classrooms, and healthcare consultations.

\section{SUGGESTIONS FOR FUTURE WORK} \label{sec:concl-future}
While OmniSign AI achieves its core architectural objectives, future research will pursue the following enhancements:
\begin{enumerate}[noitemsep, topsep=2pt]
\item \textbf{Continuous Sign Language Sentence Translation (SLT):} Extending beyond isolated word glosses by integrating spatial-temporal graph convolutional networks (ST-GCN) and connectionist temporal classification (CTC) decoders to support continuous, unsegmented sign discourse.
\item \textbf{Multilingual Indic Speech Recognition and Dialect Expansion:} Incorporating fine-tuned multilingual speech models (such as AI4Bharat's IndicWhisper) to transcribe regional Indian languages (e.g., Tamil, Hindi, Telugu) directly within the speech pipeline.
\item \textbf{Non-Manual Facial Cue and Emotion Modeling:} Enhancing MediaPipe Holistic tracking to capture eyebrow elevation, mouth morphemes, and head tilts, which serve as essential grammatical markers in sign linguistics.
\item \textbf{Edge Hardware Deployment on Embedded Systems:} Porting the ONNX and quantized INT8 models to low-power edge accelerators (such as NVIDIA Jetson Nano or Raspberry Pi 5 with Coral NPU) for portable, wearable deployment.
\end{enumerate}
\vspace{0.2cm}
\noindent\textbf{SUMMARY}\nopagebreak\par\vspace{0.1cm}
This chapter consolidated the engineering and research accomplishments of OmniSign AI, highlighting its sub-millisecond execution, 100\% test pass rate, and egalitarian bidirectional communication model. Promising avenues for future expansion include continuous sign language translation, Indic speech integration, facial expression tracking, and embedded edge deployment.
\clearpage

% ═══════════════════════════════════════════════════════════════════════════
% REFERENCES (APA Format - Strictly 2 Pages with No Lone Reference)
% ═══════════════════════════════════════════════════════════════════════════
\chapter*{REFERENCES} \label{chap:refs}
\addcontentsline{toc}{chapter}{REFERENCES}
\thispagestyle{plain}
\begin{spacing}{1.2}
\begin{enumerate}[label={[\arabic*]}, leftmargin=2.5em, itemsep=2pt]
\item Carreira, J., \& Zisserman, A. (2017). Quo vadis, action recognition? A new model and the kinetics dataset. In \textit{Proceedings of the IEEE Conference on Computer Vision and Pattern Recognition (CVPR)} (pp. 6299--6308). [https://doi.org/10.1109/CVPR.2017.502](https://doi.org/10.1109/CVPR.2017.502)

\item Chung, H. W., Hou, L., Longpre, S., Zoph, B., Tay, Y., Fedus, W., Li, E., Wang, X., Mostajabdaveh, M., Bousquet, O., \& Chen, Y. (2022). Scaling instruction-finetuned language models. \textit{arXiv preprint arXiv:2210.11416}. [https://doi.org/10.48550/arXiv.2210.11416](https://doi.org/10.48550/arXiv.2210.11416)

\item Donahue, J., Anne Hendricks, L., Guadarrama, S., Rohrbach, M., Venugopalan, S., Saenko, K., \& Darrell, T. (2015). Long-term recurrent convolutional networks for visual recognition and description. In \textit{Proceedings of the IEEE Conference on Computer Vision and Pattern Recognition (CVPR)} (pp. 2625--2634).

\item Hemg. (2023). \textit{Indian Sign Language Classification: Fine-Tuned Vision Model on ISL Alphabets}. Hugging Face Model Repository. Retrieved from [https://huggingface.co/Hemg/Indian-sign-language-classification](https://huggingface.co/Hemg/Indian-sign-language-classification)

\item Hochreiter, S., \& Schmidhuber, J. (1997). Long short-term memory. \textit{Neural Computation}, 9(8), 1735--1780. [https://doi.org/10.1162/neco.1997.9.8.1735](https://doi.org/10.1162/neco.1997.9.8.1735)

\item Koller, O., Ney, H., \& Bowden, R. (2016). Deep hand: How to train a CNN on 1 million hand images when your data is continuous and unlabelled. In \textit{Proceedings of the IEEE Conference on Computer Vision and Pattern Recognition (CVPR)} (pp. 3793--3802).

\item Li, D., Rodriguez, C., Yu, X., \& Li, H. (2020). Word-level deep sign language recognition from video: A new large-scale dataset and methods comparison. In \textit{Proceedings of the IEEE/CVF Winter Conference on Applications of Computer Vision (WACV)} (pp. 1459--1469).

\item Lugaresi, C., Tang, J., Nash, H., McClanahan, C., Uboweja, E., Hays, M., Zhang, F., Chang, C.-L., Yong, M. G., Lee, J., \& Wan, W. (2019). MediaPipe: A framework for building perception pipelines. \textit{arXiv preprint arXiv:1906.08172}. [https://doi.org/10.48550/arXiv.1906.08172](https://doi.org/10.48550/arXiv.1906.08172)

\item Ministry of Law and Justice. (2016). \textit{The Rights of Persons with Disabilities Act, 2016 (Act No. 49 of 2016)}. The Gazette of India, Extraordinary, Part II, Section 1. New Delhi: Government of India.

\item ONNX Runtime Developers. (2024). \textit{ONNX Runtime: Cross-platform, High Performance ML Inferencing and Training Accelerator}. Microsoft Corporation. Retrieved from [https://onnxruntime.ai](https://onnxruntime.ai)

\item Povey, D., Ghoshal, A., Boulianne, G., Burget, L., Glembek, O., Goel, N., Hannemann, M., Motlicek, P., Qian, Y., Schwarz, P., \& Silovsky, J. (2011). The Kaldi speech recognition toolkit. In \textit{IEEE Workshop on Automatic Speech Recognition and Understanding (ASRU)}.

\item Radford, A., Kim, J. W., Xu, T., Brockman, G., McLeavey, C., \& Sutskever, I. (2023). Robust speech recognition via large-scale weak supervision. In \textit{International Conference on Machine Learning (ICML)} (pp. 28492--28518). PMLR.

\item Rastgoo, R., Kiani, K., \& Escalera, S. (2021). Sign language recognition: A deep survey. \textit{Expert Systems with Applications}, 164, 113794. [https://doi.org/10.1016/j.eswa.2020.113794](https://doi.org/10.1016/j.eswa.2020.113794)

\item Shneiderman, B. (2020). Human-centered artificial intelligence: Reliable, safe \& trustworthy. \textit{International Journal of Human--Computer Interaction}, 36(6), 495--504. [https://doi.org/10.1080/10447318.2020.1741118](https://doi.org/10.1080/10447318.2020.1741118)

\item Starner, T., Weaver, J., \& Pentland, A. (1998). Real-time American sign language recognition using desk and wearable computer based video. \textit{IEEE Transactions on Pattern Analysis and Machine Intelligence}, 20(12), 1371--1375.

\item Stokoe, W. C. (2005). Sign language structure: An outline of the visual communication systems of the American deaf. \textit{Journal of Deaf Studies and Deaf Education}, 10(1), 3--37. [https://doi.org/10.1093/deafed/eni001](https://doi.org/10.1093/deafed/eni001)

\item Valli, C., \& Lucas, C. (2000). \textit{Linguistics of American Sign Language: An Introduction} (3rd ed.). Washington, DC: Gallaudet University Press.

\item World Health Organization (WHO). (2024). \textit{Deafness and Hearing Loss: Global Epidemiological Factsheet}. Geneva, Switzerland: World Health Organization.

\item Zhang, C. (2012). Microsoft Kinect sensor and its effect. \textit{IEEE Multimedia}, 19(2), 4--10. [https://doi.org/10.1109/MMUL.2012.24](https://doi.org/10.1109/MMUL.2012.24)

\item Zhang, F., Bazarevsky, V., Vakunov, A., Tkachenka, A., Sung, C., Chang, C.-L., \& Grundmann, M. (2020). MediaPipe Hands: On-device real-time hand tracking. \textit{arXiv preprint arXiv:2006.10214}. [https://doi.org/10.48550/arXiv.2006.10214](https://doi.org/10.48550/arXiv.2006.10214)
\end{enumerate}
\end{spacing}
\clearpage

% ═══════════════════════════════════════════════════════════════════════════
% APPENDICES
% ═══════════════════════════════════════════════════════════════════════════
\appendix
\chapter*{APPENDIX I: BILL OF MATERIALS / SYSTEM SPECIFICATIONS}
\addcontentsline{toc}{chapter}{APPENDIX I: BILL OF MATERIALS / SYSTEM SPECIFICATIONS}
\thispagestyle{plain}
\vspace{0.2cm}
\noindent Table~\ref{tab:app_bom} provides the comprehensive Bill of Materials and computational system specifications required to deploy and execute OmniSign AI in real-time environments.

\begin{table}[htbp]
\centering
\caption{System Specifications and Bill of Materials}
\label{tab:app_bom}
\small
\begin{tabularx}{\textwidth}{P{3.2cm}P{4.2cm}L}
\toprule
\textbf{Category} & \textbf{Component Specification} & \textbf{Operational Function} \\
\midrule
Host Operating System & Windows 11 / Ubuntu 22.04 LTS & Host operating environment with native audio drivers. \\
Processor (CPU) & Intel Core i5/i7 (8+ Cores), 2.4\,GHz+ & Multi-threaded landmark regression and audio capture. \\
System Memory (RAM) & 16GB DDR4 / DDR5 & In-memory ring buffer storage and model weights. \\
Graphics Card (GPU) & NVIDIA RTX 3050+ (Optional) & Optional acceleration for WLASL 3D video tensors. \\
Video Ingestion Device & USB / Built-in HD Webcam (720p/1080p) & 30 FPS visual stream capture for sign tracking. \\
Audio Ingestion Device & Standard Omnidirectional Microphone & 16\,kHz 32-bit floating-point mono audio capture. \\
Audio Playback & Integrated Stereo Speakers / Headphone & Real-time Text-to-Speech vocalization output. \\
Python Environment & Python v3.10 / v3.11 & Standard runtime hosting deep learning dependencies. \\
Skeletal Keypoint Engine & Google MediaPipe Hands v0.10+ & 21-point 3D Cartesian hand landmark extraction. \\
Landmark Inference Engine & ONNX Runtime v1.17+ & Sub-millisecond static ASL gesture classification. \\
Speech-to-Text Model & OpenAI Whisper (\texttt{whisper-base}) & Multilingual continuous ASR with RTF $< 0.2\times$. \\
NLP Grammar Engine & Google FLAN-T5-small (60M) & Local sequence-to-sequence syntactic reconstruction. \\
Web Application Stack & FastAPI, Uvicorn, Vanilla CSS & Modern glassmorphic web dashboard with WebSockets. \\
Testing Framework & Pytest v8.0+ & Automated test suite validating pipelines and buffers. \\
\bottomrule
\end{tabularx}
\end{table}
\clearpage

% ═══════════════════════════════════════════════════════════════════════════
% APPENDIX II: KEY SOURCE CODE EXTRACTS (Strictly Single Page)
% ═══════════════════════════════════════════════════════════════════════════
\chapter*{APPENDIX II: KEY SOURCE CODE EXTRACTS}
\addcontentsline{toc}{chapter}{APPENDIX II: KEY SOURCE CODE EXTRACTS}
\thispagestyle{plain}
\noindent Listing~\ref{lst:ccl_code} provides the core Python implementation of the \texttt{CommonTextLayer}, which manages multi-source stream convergence, chronological deduplication, and thread-safe listener dispatching.

\begin{lstlisting}[language=Python, caption={Common Communication Layer Implementation in Python}, label={lst:ccl_code}]
from dataclasses import dataclass, field
import threading, time
from collections import deque
from typing import Optional, Dict, Any, List, Callable

@dataclass
class TextEntry:
    text: str
    source: str               # "ISL", "ASL", "SPEECH"
    mode: str = "WORD"        # "ALPHABET", "WORD", "SPEECH"
    confidence: float = 1.0
    timestamp: float = field(default_factory=time.time)
    metadata: Dict[str, Any] = field(default_factory=dict)

class CommonTextLayer:
    """Aggregates recognition results from all sources into a unified text stream."""
    def __init__(self, max_history: int = 100, dedup_window_seconds: float = 1.5,
                 enable_deduplication: bool = True) -> None:
        self.max_entries = max_history
        self.dedup_window_sec = dedup_window_seconds
        self.enable_deduplication = enable_deduplication
        self._entries: deque[TextEntry] = deque(maxlen=self.max_entries)
        self._lock = threading.Lock()
        self._listeners: List[Callable[[TextEntry], None]] = []

    def add_entry(self, text: str, source: str = "UNKNOWN", mode: str = "WORD",
                  confidence: float = 1.0) -> Optional[TextEntry]:
        if not text or not text.strip():
            return None
        clean_text = text.strip()
        source_up, now = source.upper(), time.time()
        with self._lock:
            if self.enable_deduplication:
                for entry in reversed(self._entries):
                    if now - entry.timestamp > self.dedup_window_sec:
                        break
                    if entry.text == clean_text and entry.source == source_up:
                        return None  # Suppress duplicate token within window
            entry = TextEntry(text=clean_text, source=source_up, mode=mode.upper(),
                              confidence=confidence, timestamp=now)
            self._entries.append(entry)
        for listener in self._listeners:
            listener(entry)
        return entry
\end{lstlisting}
\clearpage

% ═══════════════════════════════════════════════════════════════════════════
% APPENDIX III: INDIVIDUAL WORK CONTRIBUTION (Strictly Single Page)
% ═══════════════════════════════════════════════════════════════════════════
\chapter*{APPENDIX III: INDIVIDUAL WORK CONTRIBUTION}
\addcontentsline{toc}{chapter}{APPENDIX III: INDIVIDUAL WORK CONTRIBUTION}
\thispagestyle{plain}
\noindent In accordance with institutional guidelines for three-candidate undergraduate project teams, Table~\ref{tab:app_contrib} details the balanced division of research, architectural design, coding, and validation responsibilities.

\begin{table}[htbp]
\centering
\caption{Individual Work Contribution Mapping}
\label{tab:app_contrib}
\small
\begin{tabularx}{\textwidth}{P{3.4cm}P{4.2cm}L}
\toprule
\textbf{Student Investigator} & \textbf{Core Module Assigned} & \textbf{Key Responsibilities \& Deliverables} \\
\midrule
\textbf{JEEVITHA S}\newline(7376222IT165) &
Visual Sign Pipelines \& Spatial Modeling &
\begin{itemize}[leftmargin=*, noitemsep, topsep=0pt]
\item Researched and integrated ISL Vision Transformer (\texttt{Hemg/ISL-Classifier}, 35 classes).
\item Developed MediaPipe 21-point hand tracking with 63-D wrist normalization.
\item Built ONNX-optimized ASL landmark classifier running at 0.086\,ms latency.
\item Integrated WLASL I3D 2,000-gloss dynamic video sequence adapter.
\end{itemize} \\
\addlinespace
\textbf{KAVIN PRASATH R}\newline(7376222IT182) &
Acoustic ASR, Stabilization, \& Convergence &
\begin{itemize}[leftmargin=*, noitemsep, topsep=0pt]
\item Integrated OpenAI Whisper ASR (\texttt{whisper-base}) on 16\,kHz PCM streams.
\item Codified 5-state gesture debouncer FSM and temporal majority stabilizer.
\item Engineered thread-safe \texttt{CommonTextLayer} with deduplication and source badges.
\item Designed and executed 40-test automated verification suite in Pytest.
\end{itemize} \\
\addlinespace
\textbf{DEEPAN RAJ M}\newline(7376222IT134) &
Neural NLP, Audio TTS, \& User Interfaces &
\begin{itemize}[leftmargin=*, noitemsep, topsep=0pt]
\item Implemented local neural grammar reconstruction using \texttt{google/flan-t5-small}.
\item Developed heuristic rule-based syntax cleaner for 100\% offline fallback.
\item Integrated real-time SAPI5 / \texttt{pyttsx3} Text-to-Speech audio engine.
\item Designed responsive glassmorphic web dashboard and OpenCV desktop HUD.
\end{itemize} \\
\bottomrule
\end{tabularx}
\end{table}
\clearpage

% ═══════════════════════════════════════════════════════════════════════════
% APPENDIX IV: PLAGIARISM VERIFICATION REPORT (Strictly Single Page)
% ═══════════════════════════════════════════════════════════════════════════
\chapter*{APPENDIX IV: PLAGIARISM VERIFICATION REPORT}
\addcontentsline{toc}{chapter}{APPENDIX IV: PLAGIARISM VERIFICATION REPORT}
\thispagestyle{plain}
\noindent The complete text of this academic project report was subjected to computerized anti-plagiarism analysis in compliance with the academic integrity guidelines mandated by Anna University and Bannari Amman Institute of Technology.
\vspace{0.3cm}

\begin{center}
\includegraphics[width=0.55\textwidth]{plagiarism-report.png}\\[4pt]
{\small\textbf{Figure IV.1:} Academic Plagiarism Verification Analysis Certificate}
\end{center}
\vspace{0.3cm}

\noindent \textbf{Summary of Plagiarism Verification:}
\begin{itemize}[noitemsep, topsep=2pt]
\item \textbf{Total Similarity Index:} 4\% (Well below the permissible statutory threshold of 15\%)
\item \textbf{Primary Source Breakdown:} Internet Sources: 2\%, Publications: 1\%, Student Papers: 1\%
\item \textbf{Verification Status:} Accepted as original research work with zero uncredited textual reproduction.
\end{itemize}
\end{document}
